BFV、BGV 和 CKKS 的参数可能需要超过机器字长的系数模数。剩余数系统（RNS，\autoref{subsec:crt-application}）使用较小且两两互素模数下的剩余类来表示这些系数，使精确模算术可以由机器字长运算实现，而不缩小总模数或改变基础安全参数。所需模数大小取决于方案、电路、精度和安全目标，没有普遍适用的固定比特长度。

CRT 使模加法与模乘法易于实现，但改变基以及实现缩放舍入还需要额外算法。本节针对 BFV、BGV 和 CKKS 介绍这些原语。常见 TFHE 实现采用可由一个机器字表示的环面系数，因此不需要这种特定的大整数分解；这一实现选择，并不是禁止其他 TFHE 变体采用 RNS 技术的定理。

本节中，$[x]_q=|x|_q$ 表示区间 $[-q/2,q/2)$ 中的唯一整数代表元（\autoref{subsec:modulo-centered}）。除非另有说明，RNS 模数均为两两互素的奇整数；使用素数模数很方便，但中国剩余定理本身并不要求它们是素数。对于奇数 $q$，最小与最大代表元分别是 $-(q-1)/2$ 和 $(q-1)/2$。必须区分这些整数代表元的界与同余关系，这对近似和精确基转换都至关重要。

\begin{tcolorbox}[
    title = \textbf{所需背景知识},    % box title
    colback = white,    % light background; tweak to taste
    colframe = black,  % frame colour
    boxrule = 0.8pt,     % line thickness
    left = 1mm, right = 1mm, top = 1mm, bottom = 1mm % inner padding
]

\begin{itemize}
\item \autoref{sec:modulo}: \nameref{sec:modulo}
\item \autoref{sec:group}: \nameref{sec:group}
\item \autoref{sec:field}: \nameref{sec:field}
\item \autoref{sec:polynomial-ring}: \nameref{sec:polynomial-ring}
\item \autoref{sec:decomp}: \nameref{sec:decomp}
\item \autoref{sec:modulus-rescaling}: \nameref{sec:modulus-rescaling}
\item \autoref{sec:chinese-remainder}: \nameref{sec:chinese-remainder}
\item \autoref{sec:polynomial-interpolation}: \nameref{sec:polynomial-interpolation}
\item \autoref{sec:ntt}: \nameref{sec:ntt}
\item \autoref{sec:lattice}: \nameref{sec:lattice}
\item \autoref{sec:rlwe}: \nameref{sec:rlwe}
\item \autoref{sec:glwe}: \nameref{sec:glwe}
\item \autoref{sec:glwe-add-cipher}: \nameref{sec:glwe-add-cipher}
\item \autoref{sec:glwe-add-plain}: \nameref{sec:glwe-add-plain}
\item \autoref{sec:glwe-mult-plain}: \nameref{sec:glwe-mult-plain}
\item \autoref{subsec:modulus-switch-rlwe}: \nameref{subsec:modulus-switch-rlwe}
\item \autoref{sec:glwe-key-switching}: \nameref{sec:glwe-key-switching}
\item \autoref{sec:bfv}: \nameref{sec:bfv}
\item \autoref{sec:ckks}: \nameref{sec:ckks}
\item \autoref{sec:bgv}: \nameref{sec:bgv}
\end{itemize}
\end{tcolorbox}

\clearpage

\subsection{快速基转换：\textsf{FastBConv}}
\label{subsec:rns-fastbconv}

\noindent \textbf{- 参考文献：} \href{https://eprint.iacr.org/2016/510}{A Full RNS Variant of FV-like Somewhat Homomorphic Encryption Schemes}~\cite{rns-bfv}；\href{https://eprint.iacr.org/2022/657}{BASALISC: Programmable Hardware Accelerator for BGV FHE}~\cite{rns-bfv2}。

令 $q=\prod_{i=1}^kq_i$、$b=\prod_{j=1}^lb_j$，所有基模数两两互素。输入剩余类 $x_i=[x]_{q_i}$ 确定中心化整数 $x=[x]_q$。精确基转换应输出 $([x]_{b_j})_{j=1}^l$。快速基转换省略一次大模数约减，而在每个输出坐标使用同一个未经约减的 CRT 和：
\[
y_i=q/q_i,\qquad z_i=y_i^{-1}\bmod q_i,\qquad
T=\sum_{i=1}^k[x_iz_i]_{q_i}y_i.
\]
\begin{tcolorbox}[title={\textbf{\tboxlabel{\ref*{subsec:rns-fastbconv}} 快速基转换：\textsf{FastBConv}}}]
\[
\textsf{FastBConv}(x,q,b)=([T]_{b_1},\ldots,[T]_{b_l}).
\]
存在整数 $u$ 使 $T=x+uq$。对于上述奇模数中心化约定，
\[
|T|<kq/2,\quad |x|<q/2,\quad |u|<(k+1)/2,
\quad |u|\leq\lfloor k/2\rfloor.
\]
因此输出表示 $x+uq\pmod b$，不一定表示 $x\pmod b$。记该界为 $U_k=\lfloor k/2\rfloor$。
\end{tcolorbox}

\begin{myproof}
每个加项的绝对值小于 $q/2$，且 CRT 同余关系给出 $T\equiv x\pmod q$。对 $u=(T-x)/q$ 的界必须同时考虑 $T$ 与 $x$；只约束 $T$，不能将同一区间直接用作 $uq$ 的界。所有输出剩余类都由共同的 $T$ 计算，因此误差是在整个目标基上一致的倍数 $uq$。其他剩余类约定会得到不同的界，包括~\cite{rns-bfv} 采用的无符号约定。
\end{myproof}

这是基转换的近似实现，而不是中国剩余定理失效。后续算法必须消除该 $q$ 的倍数，将其通过除法转化为受控误差，或明确将其计入正确性界。

\subsection{小蒙哥马利约减算法：\textsf{SmallMont}}
\label{subsec:rns-smallmont}

Small Montgomery 约减限制快速基转换引入的 $q$ 的倍数，一般不会使基转换变为精确操作。下面的中心化剩余类变体明确列出所需条件；原始无符号剩余类算法及其参数界见~\cite{rns-bfv} 的算法 2 和引理 4。

令 $x=[x]_q$，选取与 $q$、$b$ 均互素的奇数辅助模数 $b_\alpha$，并计算 $[b_\alpha x]_q$ 的快速 CRT 提升：
\[
T=[b_\alpha x]_q+uq=b_\alpha x+(v+u)q,
\qquad |u|\leq U_k.
\]
输入由同一个 $T$ 在各个 $b_j$ 及 $b_\alpha$ 下的剩余类组成。
\begin{tcolorbox}[title={\textbf{\tboxlabel{\ref*{subsec:rns-smallmont}} Small Montgomery 约减：\textsf{SmallMont}}}]
\begin{enumerate}
\item 使用给定的 $b_\alpha$ 剩余类，计算 $a=[Tq^{-1}]_{b_\alpha}$。
\item 对 $j=1,\ldots,l$，输出 $r_j=[(T_{b_j}-qa)b_\alpha^{-1}]_{b_j}$。
\end{enumerate}
结果表示 $r=(T-qa)/b_\alpha=x+u'q\pmod b$。若 $b_\alpha>U_k+1/2$，则 $u'\in\{-1,0,1\}$。
\end{tcolorbox}

\begin{myproof}
因为 $a\equiv v+u\pmod{b_\alpha}$，商 $u'=(v+u-a)/b_\alpha$ 是整数。中心化代表元给出 $|v|<(b_\alpha+1)/2$ 和 $|a|<b_\alpha/2$，因此
\[
|u'|<1+\frac{U_k+1/2}{b_\alpha}<2.
\]
这就证明了所述的三种可能整数值。分子 $T-qa$ 被 $b_\alpha$ 整除，可以在每个 $b_j$ 下通过乘以 $b_\alpha^{-1}\bmod b_j$ 计算其商的剩余类。
\end{myproof}

计算过程中必须始终保留辅助剩余类。先重构 $[T]_b$，再将该值对 $b_\alpha$ 取模，通常会丢失原来的 $b_\alpha$ 剩余类。此外，$v$ 可能随 $b_\alpha$ 线性增长；声称其增长更慢并不能证明该界。

\subsubsection{使用 \textsf{SmallMont} 改进 \textsf{FastBConv}}
\label{subsubsec:rns-smallmont-fastbconv}

满足 \autoref{subsec:rns-smallmont} 的参数条件时，
\[
\textsf{SmallMont}\bigl(\textsf{FastBConv}([b_\alpha x]_q,q,bb_\alpha),bb_\alpha,b_\alpha,q\bigr)
=[x+u'q]_b,\qquad u'\in\{-1,0,1\}.
\]
预先乘以 $b_\alpha$ 是必要步骤。结果仍是中心化 $x$ 的近似提升，与 \autoref{subsec:rns-fastbconvex} 的 \textsf{FastBConvEx} 不同。


\subsection{基于 RNS 的 ModRaise：\textsf{ModRaise\textsubscript{RNS}}}
\label{subsec:rns-modraise}

\noindent \textbf{- 参考文献：} \href{https://eprint.iacr.org/2018/931}{A Full RNS Variant of Approximate Homomorphic Encryption}~\cite{rns-ckks}，第 3.2 节。

从 $q$ 到 $qb$ 的基扩展保留模各个 $q_i$ 的剩余类，并附加模各个 $b_j$ 的剩余类。使用快速基转换时，它选择一个近似整数提升：
\begin{tcolorbox}[title={\textbf{\tboxlabel{\ref*{subsec:rns-modraise}} \textsf{ModRaise\textsubscript{RNS}}}}]
\[
\textsf{ModRaise}_{\mathrm{RNS}}(x,q,qb)
=(x_1,\ldots,x_k,\textsf{FastBConv}(x,q,b)).
\]
组合后的向量表示 $x+uq\pmod{qb}$，其中 $|u|\leq U_k$。采用 \autoref{subsec:rns-smallmont} 中的 Montgomery 预处理，可以将其限制为满足 $|u'|\leq1$ 的适当提升 $x+u'q$。
\end{tcolorbox}

这是整数代表元的选择，并不是从 $\mathbb Z_q$ 到 $\mathbb Z_{qb}$ 的典范环嵌入。若逐系数作用于密文分量，将新代表元写为 $A+qU_A$ 和 $B+qU_B$，则其相位改变
\[
q(U_B+U_AS).
\]
因此，分量界 $\|U_A\|_\infty,\|U_B\|_\infty\leq U_k$ 给出的相位界是 $qU_k(1+\|S\|_1)$，而不只是 $qU_k$。整数相位 $B+AS$ 中还可能已经含有 $q$ 的倍数。这些项必须由后续乘法或自举算法处理，不能视为独立的新 RLWE 误差。

\subsection{基于 RNS 的 ModDrop：\textsf{ModDrop\textsubscript{RNS}}}
\label{subsec:rns-moddrop}

\textsf{ModDrop\textsubscript{RNS}} 从模 $qb$ 的表示中删除对应于 $b$ 的各因子的剩余类：
\begin{tcolorbox}[title={\textbf{\tboxlabel{\ref*{subsec:rns-moddrop}} \textsf{ModDrop\textsubscript{RNS}}}}]
\[
(x_{q_1},\ldots,x_{q_k},x_{b_1},\ldots,x_{b_l})
\longmapsto(x_{q_1},\ldots,x_{q_k}).
\]
这是精确的模 $q$ 约减，不进行缩放或舍入。
\end{tcolorbox}

对于 CKKS 或 BGV，只有原中心化整数相位的每个系数都已经位于 $[-q/2,q/2)$ 内，才能保持同一相位。此时 CKKS 保持尺度不变，BGV 保持明文模数不变。否则，相位回绕可能改变解码消息。对于 BFV，即使相位不变，也通常不能保持消息，因为标准尺度从 $\lfloor qb/t\rfloor$ 变为 $\lfloor q/t\rfloor$。例如，初始模数为 $105$、$t=7$，无噪声相位 $15$ 编码 $m=1$；降至模数 $35$ 后相位仍为 $15$，但标准 BFV 解码得到 $3$。需要保持消息时，应采用 BFV 专用的缩放舍入操作。

\subsection{基于 RNS 的模数切换：\textsf{ModSwitch\textsubscript{RNS}}}
\label{subsec:rns-modswitch}

令整数 $x$ 以 $q$ 与 $b$ 的组合基给出。对于奇数 $b$，定义 $r=[x]_b$，则 $(x-r)/b=\lfloor x/b\rceil$。下面的操作近似除以 $b$，并删除其基因子：
\begin{tcolorbox}[title={\textbf{\tboxlabel{\ref*{subsec:rns-modswitch}} \textsf{ModSwitch\textsubscript{RNS}}}}]
\begin{enumerate}
\item 使用 \textsf{FastBConv} 将 $r$ 的剩余类从基 $b$ 转换到基 $q$。其共同整数提升是 $r+ub$，其中 $|u|\leq U_l$。
\item 对每个 $i=1,\ldots,k$，输出
\[
y_i=[b^{-1}(x_{q_i}-(r+ub)_{q_i})]_{q_i}.
\]
\end{enumerate}
结果表示
\[
y=\left\lfloor\frac{x}{b}\right\rceil-u\pmod q,\qquad |u|\leq U_l.
\]
\end{tcolorbox}

代入 $x-r=b\lfloor x/b\rceil$ 即可得到该式。此同余关系不要求 $b\ll q$。只删除一个奇模数时，$l=1$ 且 $u=0$，结果就是精确的最近整数除法。删除多个因子时，必须将额外误差 $-u$ 计入密文相位界。作用于 $(A,B)$ 后，总扰动还包含分量误差与 $S$ 的乘积，不能只用单个系数的 $u$ 作为其界。

CKKS 可将其用作重缩放原语，同时跟踪实际除以 $b$ 后的尺度和全部舍入误差。BFV 和 BGV 需要各自的明文保持规则。特别地，普通的缩放舍入并不是 \autoref{subsec:rns-modswitch-bgv} 中的 BGV 模数切换。

\subsubsection{\textsf{ModSwitch\textsubscript{RNS}}、\textsf{ModRaise\textsubscript{RNS}} 与 \textsf{ModDrop\textsubscript{RNS}} 的比较}

\textsf{ModRaise\textsubscript{RNS}} 选择进入更大组合基的提升，\textsf{ModDrop\textsubscript{RNS}} 精确约减到更小的子基，\textsf{ModSwitch\textsubscript{RNS}} 则还会对所表示整数作近似除法。这是不同的操作。它们能否保持密文的明文，取决于具体方案、尺度或明文模数，以及完整的中心化相位界。




\subsection{基于 RNS 的解密}
\label{subsec:rns-dec}

本小节介绍如何高效解密采用 RNS 表示的 BFV、CKKS 与 BGV 密文。


\subsubsection{BFV 解密：$\textsf{Dec}_{\textsf{RNS}}^{\textsf{BFV}}$}
\label{subsubsec:rns-dec-bfv}

采用本书的密文约定 $B=-AS+\Delta M+E\pmod q$。对每个中心化相位系数 $x=[B+AS]_q$，标准 BFV 解码操作是 $[\lfloor tx/q\rceil]_t$；除以 $\Delta=\lfloor q/t\rfloor$ 并不是代数上相同的运算。下面的冗余模数纠正在满足舍入余量时实现缩放舍入，参见~\cite{rns-bfv} 第 3 节。

令 $q=\prod_{i=1}^kq_i$ 为奇数、$\gcd(q,t)=1$，并选择与 $qt$ 互素的奇数 $\gamma$。记
\[
m_0=\left\lfloor\frac{tx}{q}\right\rceil,\quad
\delta=\frac{tx}{q}-m_0,\quad v=[t\gamma x]_q.
\]
$v$ 的快速 CRT 提升为 $v+uq$，其中 $|u|\leq U_k$。计算其模 $\gamma$ 和 $t$ 的剩余类，并分别乘以对应模数下的 $-q^{-1}$，得到 $(y_\gamma,y_t)$。于是
\[
y\equiv-q^{-1}(v+uq)
\equiv\left\lfloor\frac{t\gamma x}{q}\right\rceil-u
=\gamma m_0+\eta\pmod{\gamma t},
\qquad \eta=\lfloor\gamma\delta\rceil-u.
\]
若 $|\eta|<\gamma/2$，其中心化剩余类恰为 $y_\gamma=\eta$，因此
\[
[m_0]_t=[(y_t-y_\gamma)\gamma^{-1}]_t.
\]
一个充分条件是 $\gamma|\delta|+1/2+U_k<\gamma/2$。

\begin{tcolorbox}[title={\textbf{\tboxlabel{\ref*{subsubsec:rns-dec-bfv}} $\textsf{Dec}_{\textsf{RNS}}^{\textsf{BFV}}$}}]
\begin{enumerate}
\item 计算相位剩余类 $x_i=[B+AS]_{q_i}$。
\item 选择满足互素条件和舍入余量界的 $\gamma$，将 $[t\gamma x]_q$ 快速转换到两个模数 $\gamma,t$。
\item 将两个剩余类分别乘以对应模数下的 $-q^{-1}$，得到 $(y_\gamma,y_t)$，其中 $y_\gamma$ 取中心化代表元。
\item 输出 $[(y_t-y_\gamma)\gamma^{-1}]_t$。
\end{enumerate}
\end{tcolorbox}

BFV 正确解密还要求 $[m_0]_t$ 等于预期明文。仅选择任意互素的 $\gamma$ 并不足够，而且纠正阈值是 $\gamma/2$，不是 $\gamma$。任意接近舍入边界的值需要相应更强的精度或纠正策略。这不是一个适用于任意输入与模数的无条件精确舍入原语。\cite{rns-bfv-hps} 中其他缩放舍入构造也规定了各自的精度和误差条件。


\subsubsection{CKKS 与 BGV 解密}
\label{subsubsec:rns-dec-ckks-bgv}

在当前 RNS 基中计算 $[B+AS]_q$。CKKS 按照已跟踪的尺度解码其中心化整数代表元，BGV 则将该代表元对明文模数取模。CRT 重构是一种有效实现。只有整个中心化相位都位于较小模数的一半以内，才能在解码前直接降到单个小模数。当前层级的有效密文不一定满足这个更强的界。必要时，应使用符合方案的重缩放或模数切换并分析误差，不能假设单纯 \textsf{ModDrop} 到最低层始终保持消息。

\subsection{BGV 的 RNS 模数切换：$\textsf{ModSwitch}_{\textsf{RNS}}^{\textsf{BGV}}$}
\label{subsec:rns-modswitch-bgv}

BGV 模数切换在降低密文模数的同时，必须保持明文剩余类。\autoref{subsec:rns-modswitch} 中的普通舍入原语不保证这一条件。也不能通过令 $t=\hat q$ 调用 $\textsf{Dec}_{\textsf{RNS}}^{\textsf{BFV}}$ 来获得无条件的舍入操作：若 $\hat q$ 整除 $q_l$，所需逆元 $q_l^{-1}\bmod\hat q$ 根本不存在。

下面给出一种常用模数链步骤的直接 RNS 实现。令 $q=Qw$，满足 $\gcd(w,Q)=\gcd(w,t)=1$ 和 $w\equiv1\pmod t$。对任一密文分量的每个中心化系数 $a$，定义
\[
d=t[t^{-1}a]_w,\qquad a'=\frac{a-d}{w}\pmod Q.
\]
由于 $d\equiv a\pmod w$ 且 $d\equiv0\pmod t$，这是整数除法，而且在对 $Q$ 约减前满足 $a'\equiv a\pmod t$。在每个保留模数 $q_i\mid Q$ 下，通过下式计算：
\[
a'_i=[w^{-1}(a_i-d)]_{q_i}.
\]
若 $w$ 含有多个 RNS 因子，必须重构中心化剩余类 $[t^{-1}a]_w$，或采用精确且具有已证明界限的转换过程；单因子步骤不需要大整数重构。

令输入相位为 $M+tE$，并使用 \autoref{subsec:bgv-modulus-switch} 中同一个溢出多项式。约减前的输出相位为
\[
Y=\frac{M+tE}{w}-\frac{d_B+d_AS}{w}.
\]
由于 $q=Qw\equiv Q\pmod t$，该相位在模 $t$ 下与 $M$ 同余。一个充分的不回绕条件是
\[
\frac{\|M+tE\|_\infty}{w}+\frac{t}{2}(1+\|S\|_1)<\frac Q2.
\]
在此条件下，输出是相同明文的有效 BGV 密文。更一般的模数链可以跟踪明文纠正因子；这里的同余假设是明确声明的特例，并不是对所有 BGV 的限制。


\subsection{精确快速基转换：\textsf{FastBConvEx}}
\label{subsec:rns-fastbconvex}

与 \textsf{SmallMont} 不同，带冗余模数的精确基转换在满足明确界限时会消除全部 CRT 歧义。这是~\cite{rns-bfv} 第 4.4 节引理 6 使用的 Shenoy--Kumaresan 型纠正。令 $b=\prod_{j=1}^lb_j$，假设已知整数 $x$ 在每个 $b_j$ 以及额外互素奇模数 $b_\alpha$ 下的剩余类。写成
\[
\widehat x=[x]_b=x+\mu b,\qquad
T=\widehat x+ub=x+(\mu+u)b,
\]
其中 $T$ 是两次快速转换共同使用的未经约减的 CRT 和，且 $|u|\leq U_l$。要求 $|\mu+u|<b_\alpha/2$；若已知 $|\mu|\leq\lambda$，可以采用充分条件 $b_\alpha>2(\lambda+U_l)$。

\begin{tcolorbox}[title={\textbf{\tboxlabel{\ref*{subsec:rns-fastbconvex}} 精确快速基转换：\textsf{FastBConvEx}}}]
\begin{enumerate}
\item 保留 $x$ 的基 $b$ 剩余类，并单独保留给定剩余类 $x_\alpha=[x]_{b_\alpha}$。
\item 计算 $\gamma=[(\textsf{FastBConv}(\widehat x,b,b_\alpha)-x_\alpha)b^{-1}]_{b_\alpha}$。
\item 在目标基的各坐标中输出 $[\textsf{FastBConv}(\widehat x,b,q)-\gamma b]_q$。
\end{enumerate}
\end{tcolorbox}

\begin{myproof}
两次转换都使用 $T$。因此在所述界限下，$\gamma=[\mu+u]_{b_\alpha}=\mu+u$，从而 $T-\gamma b=x$，输出恰为 $[x]_q$。这是代数恒等式，不会遗留新的近似误差。若不满足该界，纠正结果可能相差 $bb_\alpha$ 的倍数，而该倍数在模 $q$ 下不一定为零。
\end{myproof}

\para{为何需要中心化（即有符号）剩余类表示：} 必须将纠正项 $\gamma$ 解释为有界有符号整数 $\mu+u$。当 $|\mu+u|<b_\alpha/2$ 时，中心化约定提供了这种解释。无符号实现也可以使用，但必须显式恢复符号；直接把无符号剩余类当作该整数是不正确的。所需条件是可量化的范围界，而不是未经量化的“$x$ 远小于模数”。


\subsection{乘法分解：\textsf{DecompMult\textsubscript{RNS}}}
\label{subsec:rns-decompmult}

RNS 为密钥切换和重线性化提供了一个整数 gadget。令 $q=\prod_iq_i$，定义
\[
g_i=\frac q{q_i}\left[\left(\frac q{q_i}\right)^{-1}\right]_{q_i},\qquad
A_i=[A]_{q_i},\qquad A\equiv\sum_iA_ig_i\pmod q.
\]
这里将每个 $A_i$ 逐系数提升为小整数多项式。将其乘以 $R_q$ 中的密文时，必须在 RNS 基的\emph{每一个}模数下约减同一组小整数。如果只保留第 $i$ 个坐标、其余填零，就会得到不同的乘数和误差界。

\begin{tcolorbox}[title={\textbf{\tboxlabel{\ref*{subsec:rns-decompmult}} \textsf{DecompMult\textsubscript{RNS}}}}]
从 $S$ 切换到 $S'$ 时，生成满足下式的求值密钥 $(a_i,b_i)$：
\[
b_i+a_iS'\equiv g_iS+E_i\pmod q.
\]
于是 $\sum_iA_i(a_i,b_i)$ 在模 $q$ 下具有相位 $AS+\sum_iA_iE_i$。加上 $(0,B)$ 即完成密钥切换。在 BGV 中应以 $tE_i$ 替代 $E_i$，使新增误差被明文模数整除。

重线性化时，将密钥中的 $g_iS$ 替换为 $g_iS^2$，分解 $D_2$，再按照本书的 $(A,B)$ 顺序，将结果加到 $(D_1,D_0)$ 上。所得相位为
\[
D_0+D_1S+D_2S^2+\sum_iD_{2,i}E_i\pmod q.
\]
\end{tcolorbox}

对于负循环多项式，$\|A_iE_i\|_\infty\leq n\|A_i\|_\infty\|E_i\|_\infty$，因此新增误差的一个明确界为 $n\sum_i(q_i/2)\|E_i\|_\infty$。也可以把模逆元放入数位 $[A_i(q/q_i)^{-1}]_{q_i}$，并采用权重 $q/q_i$，如~\cite{rns-bfv} 第 4.5 节所示。两种选择都有效，但不能混用它们的数位与密钥。这些恒等式中的加密符号表示相位关系，而不是重新随机采样的密文分布相等。


\subsection{将 RNS 技术应用于 FHE 运算}
\label{subsec:rns-application}

本小节说明如何使用此前介绍的 RNS 基本操作，处理 BFV、CKKS 与 BGV 中采用 RNS 表示的密文运算。


\subsubsection{多项式加法与乘法}
\label{subsubsec:rns-application-basic}

RNS 分解的是系数模数，因此标量加法与标量乘法分别在每个 $q_i$ 下进行（\autoref{subsec:crt-application}）。多项式加法逐系数进行，但多项式乘法仍需要模 $X^n+1$ 的负循环卷积。只有在模数支持适当 NTT 并完成变换后，乘法才变成逐点乘法。下面的标量恒等式描述了这些多项式运算内部使用的系数算术。

$P^{\langle 1 \rangle} = \sum\limits_{a=0}^{n-1} c^{\langle 1 \rangle}_a \cdot X^a \in \mathcal{R}_{\langle n, q\rangle}$

$P^{\langle 2 \rangle} = \sum\limits_{b=0}^{n-1} c^{\langle 2 \rangle}_b \cdot X^b \in \mathcal{R}_{\langle n, q\rangle}$

$ $

在 FHE 方案的 RNS 变体中，将每个多项式系数表示为如下 RNS 剩余向量：

$c^{\langle 1 \rangle}_a = (c^{\langle 1 \rangle}_{a,1}, c^{\langle 1 \rangle}_{a,2}, \cdots, c^{\langle 1 \rangle}_{a,k}) \in \mathbb{Z}_{q_1} \times \mathbb{Z}_{q_2} \times \cdots \times \mathbb{Z}_{q_k}$ \textcolor{red}{ $\rhd$ 对于 $a \in [0, n - 1]$}

$c^{\langle 2 \rangle}_b = (c^{\langle 2 \rangle}_{b,1}, c^{\langle 2 \rangle}_{b,2}, \cdots, c^{\langle 2 \rangle}_{b,k}) \in \mathbb{Z}_{q_1} \times \mathbb{Z}_{q_2} \times \cdots \times \mathbb{Z}_{q_k}$ \textcolor{red}{ $\rhd$ 对于 $b \in [0, n - 1]$}

$ $

在上述 RNS 设置下，对两个多项式执行加法或乘法时，每次系数相加都通过两个 RNS 剩余向量的逐元素加法计算：

$c^{\langle 1 \rangle}_a + c^{\langle 2 \rangle}_b \equiv \sum\limits_{i=1}^k (c^{\langle 1 \rangle}_{a,i} + c^{\langle 2 \rangle}_{b,i}) y_i z_i \bmod q $ \textcolor{red}{ $\rhd$ 其中 $y_i = \dfrac{q}{q_i}$、$z_i = |y_i^{-1}|_{q_i}$}

$\Longleftrightarrow (c^{\langle 1 \rangle}_{a,1} + c^{\langle 2 \rangle}_{b,1}, \text{ } c^{\langle 1 \rangle}_{a,2} + c^{\langle 2 \rangle}_{b,2}, \cdots, \text{ } c^{\langle 1 \rangle}_{a,k} + c^{\langle 2 \rangle}_{b,k}) \in \mathbb{Z}_{q_1} \times \mathbb{Z}_{q_2} \times \cdots \times \mathbb{Z}_{q_k}$

$ $

类似地，每次系数相乘都通过两个 RNS 剩余向量的逐元素乘法计算：

$c^{\langle 1 \rangle}_a \cdot c^{\langle 2 \rangle}_b \equiv \sum\limits_{i=1}^k (c^{\langle 1 \rangle}_{a,i} \cdot c^{\langle 2 \rangle}_{b,i}) y_i z_i \bmod q $

$\Longleftrightarrow (c^{\langle 1 \rangle}_{a,1} \cdot c^{\langle 2 \rangle}_{b,1}, \text{ } c^{\langle 1 \rangle}_{a,2} \cdot c^{\langle 2 \rangle}_{b,2}, \cdots, \text{ } c^{\langle 1 \rangle}_{a,k} \cdot c^{\langle 2 \rangle}_{b,k}) \in \mathbb{Z}_{q_1} \times \mathbb{Z}_{q_2} \times \cdots \times \mathbb{Z}_{q_k}$

$ $

利用上述同构，可以基于较小的 RNS 剩余（例如每个 30 位），高效地对大多项式系数（例如 1000 位）执行密文与明文相加、密文与密文相加以及密文与明文相乘。


\subsubsection{密钥切换}
\label{subsubsec:rns-application-key-switching}

使用 \autoref{subsec:rns-decompmult} 的求值密钥和提升数位，计算
\[
(A',B')=(0,B)+\textsf{DecompMult}_{\mathrm{RNS}}(A,S_{\langle S',\mathrm{RNS}\rangle}).
\]
若输入相位为 $B+AS$，则输出相位在模 $q$ 下为 $B+AS+E_{\mathrm{ks}}$，其中 $E_{\mathrm{ks}}=\sum_iA_iE_i$。BFV 和 CKKS 保持当前消息尺度，同时增加该误差。BGV 采用误差被 $t$ 整除的求值密钥，使相位仍在模 $t$ 下与同一消息同余。全部运算均可在当前 RNS 基中执行，但相位界仍取决于具体方案。


\subsubsection{输入槽位旋转}
\label{subsubsec:rns-rotation}

对于 $J(h)=5^h\bmod2n$，对两个密文分量施加自同构 $X\mapsto X^{J(h)}$。将指数模 $2n$ 约减，并使用 $X^{n+i}=-X^i$。在系数表示中，这相当于置换整个 RNS 剩余类向量，并在所有模数下对相应向量取负。该步骤本身不引入近似误差。

所得相位是原相位在 $X^{J(h)}$ 处的代入结果，秘密密钥变为 $S(X^{J(h)})$。随后使用 \autoref{subsubsec:rns-application-key-switching} 的 RNS 密钥切换，返回 $S(X)$。此步骤会增加密钥切换误差，BGV 的误差含有因子 $t$。产生的槽位置换取决于根的排列：CKKS 具有 $n/2$ 个独立复数槽位，而完全分裂的 BFV/BGV 打包在 $5$ 的幂作用下具有两个各含 $n/2$ 个标量槽位的轨道。RNS 不改变该置换。

\subsubsection{BFV 的密文与密文相乘}
\label{subsubsec:rns-cipher-mult-bfv}

\cite{rns-bfv} 第 4.2--4.5 节的 BEHZ RNS 乘法依次执行基扩展、张量乘法、近似缩放舍入、精确转换回原基，以及重线性化。辅助基和误差界是算法的一部分。HPS 构造~\cite{rns-bfv-hps} 第 2.2--2.4 节及第 4.5 节使用不同的转换与缩放技术，不能混淆这些变体。

下面给出本节原语的中心化剩余类版本。令 $b=\prod_jb_j$，使用冗余模数 $b_\alpha$ 进行末尾的精确转换，并使用\emph{另一个不同的}临时 Montgomery 模数 $\widetilde m$ 控制初始提升。所有基模数都必须满足所需互素条件。
\begin{enumerate}
\item \textbf{基扩展。} 对每个分量系数 $x$，将 $[\widetilde m x]_q$ 快速转换到组合辅助基 $bb_\alpha\widetilde m$，再以 $\widetilde m$ 为临时模数执行 \textsf{SmallMont}。这样得到模 $bb_\alpha$ 下的提升 $\widehat x=x+uq$，其中 $|u|\leq1$，同时保留原基 $q$ 的剩余类。辅助剩余类 $b_\alpha$ 留作后续精确转换，不能像 Montgomery 模数那样被删除。
\item \textbf{张量乘法。} 在组合基 $qbb_\alpha$ 中计算负循环整数多项式乘积
\[
D_0=\widehat B_1\widehat B_2,\quad
D_1=\widehat A_1\widehat B_2+\widehat A_2\widehat B_1,\quad
D_2=\widehat A_1\widehat A_2.
\]
选择辅助范围，使预期整数 $tD_j$ 能够无回绕表示，例如验证 $\|tD_j\|_\infty<qbb_\alpha/2$。系数卷积、$n$、$t$ 和分量提升的界都会影响这一条件；仅有 $b>\Delta$ 并不足够。
\item \textbf{近似缩放舍入。} 将各个 $D_j$ 乘以 $t$，并执行删除基 $q$ 的 \textsf{ModSwitch\textsubscript{RNS}}。输出在基 $bb_\alpha$ 中表示
\[
Z_j=\left\lfloor\frac{tD_j}{q}\right\rceil-u_j
=\frac{tD_j}{q}+\rho_j,\quad
\|u_j\|_\infty\leq U_k,\quad
\|\rho_j\|_\infty\leq U_k+\frac12.
\]
这是近似缩放舍入，不是精确舍入。
\item \textbf{精确返回基 $q$。} 对每个 $Z_j$，使用保留的 $b_\alpha$ 剩余类执行 \textsf{FastBConvEx}。每个系数都必须满足 \autoref{subsec:rns-fastbconvex} 中的界 $|\mu+u|<b_\alpha/2$。对 $tD_j$ 的辅助范围界，不能代替这一独立的纠正界。
\item \textbf{重线性化。} 按 $(A,B)$ 顺序计算 $(Z_1,Z_0)+\textsf{DecompMult}_{\mathrm{RNS}}(Z_2,S^2_{\langle S,\mathrm{RNS}\rangle})$，此步骤增加求值密钥误差 $E_{\mathrm{ks}}$。
\end{enumerate}

令 $H_i=\widehat B_i+\widehat A_iS$ 表示提升后的整数相位。输出相位满足
\[
F\equiv\frac tqH_1H_2+\rho_0+\rho_1S+\rho_2S^2+E_{\mathrm{ks}}\pmod q.
\]
不能将三个分量的舍入替换为对 $H_1H_2$ 的单次舍入。误差界必须包含依赖秘密密钥的项、输入误差、提升产生的倍数，以及 $t\Delta^2/q$ 与 $\Delta$ 的差异。例如，将中心化输出相位写为 $\Delta M_{\mathrm{out}}+E_{\mathrm{out}}$，一个充分的解码条件是
\[
\|\Delta M_{\mathrm{out}}+E_{\mathrm{out}}\|_\infty<q/2,\qquad
\left\|\frac tqE_{\mathrm{out}}+\left(\frac{t\Delta}{q}-1\right)M_{\mathrm{out}}\right\|_\infty<1/2.
\]
其中 $M_{\mathrm{out}}=[M_1M_2]_t$。这些是绝对误差条件；仅将误差与消息大小比较并不正确，消息系数为零时尤其如此。相应的完整参数分析见 \autoref{subsubsec:bfv-mult-cipher-summary} 和所引论文。


\subsubsection{CKKS 的密文与密文相乘}
\label{subsubsec:rns-cipher-mult-ckks}

直接在当前基中计算张量积，使用适当的 RNS 求值密钥重线性化，并在尺度与噪声预算要求时重缩放（\autoref{subsubsec:ckks-mult-cipher-summary}）。在使用近似解释前，乘积相位必须适合当前模数；仅凭层级编号不能保证这一点。

若输入尺度为 $\Delta_1,\Delta_2$，删除的模数为 $w$，则新尺度是 $\Delta_1\Delta_2/w$。单素数重缩放能精确实现分量的最近整数除法，但仍产生相位误差 $R_B+R_AS$。多因子的近似转换还可能引入其他有界误差。BFV 专用的因子 $t$、用于张量积的大范围辅助基，以及精确返回固定 BFV 模数的最后步骤，不属于这一 CKKS 流程。不过，混合密钥切换仍可能使用自己的辅助基。


\subsubsection{BGV 的密文与密文相乘}
\label{subsubsec:rns-cipher-mult-bgv}

张量乘法与 RNS 重线性化得到模 $t$ 下的同一明文乘积，误差仍被 $t$ 整除（\autoref{subsec:bgv-mult-cipher}）。求值密钥误差也必须是 $t$ 的倍数。如需模数切换，应使用 \autoref{subsec:rns-modswitch-bgv} 的 $\textsf{ModSwitch}_{\mathrm{RNS}}^{\mathrm{BGV}}$，或显式跟踪纠正因子的更一般同余保持变体，不能用普通 CKKS 重缩放代替。因子 $t$ 保持固定；正确性和有效的噪声降低取决于完整的输出相位界。


\subsubsection{BFV 自举}
\label{subsubsec:rns-bfv-bootstrapping}

RNS 是密文系数模数的表示方式，并不禁止明文模数为 $p^\varepsilon$。素数幂也可以作为一个 CRT 模数，只要它与其他基模数互素。无效的只是把重复因子 $p,p,\ldots,p$ 当作不同且互素的 CRT 坐标。

\cite{cryptoeprint:2022/1363} 第 3.2--3.4 节的 BFV 自举构造采用公开预处理
\[
c_i'=\left[\left\lfloor\frac{p^\varepsilon}{q}c_i\right\rceil\right]_{p^\varepsilon},
\]
随后使用自举密钥同态计算 $[c_0'+c_1'S]_{p^\varepsilon}$。这里 $(c_0,c_1)=(B,A)$，该同态计算的明文模数为 $p^\varepsilon$。线性变换与数位提取实现将明文模数降至 $p^r$ 的舍入解密映射，之后施加逆变换。方案层面的讨论见 \autoref{subsubsec:bfv-bootstrapping-summary}。

这种系数到槽位构造要求明文素数与分圆指标 $\mu$ 满足 $\gcd(p,\mu)=1$（所引论文第 2.3 节）。对于 $\mu=2n$ 且 $n$ 为二次幂的情形，应使用奇数 $p$。二进制数位提取规则本身不能补足该环所缺少的明文环 CRT 分解。

公开预处理可以从 RNS 重构密文系数并精确舍入，也可以使用精度足够且经过证明的 RNS 缩放舍入例程。此后，同态电路的密文算术可以为工作密文模数使用另一套 RNS 基。若采用近似基转换，必须先把其误差计入舍入解密界，才能声称正确性。仅有 BEHZ/HPS 乘法例程并不能证明新的自举构造，而关于因子不互素的错误论证也不能成为用未经证明的基扩展与重新解释链替换该预处理的理由。

对于理想化输入相位 $(q/p^r)M+E$，分量舍入贡献 $R_B+R_AS$，舍入解密误差包含 $(p^\varepsilon/q)E+R_B+R_AS$。若实际加密尺度为 $\lfloor q/p^r\rfloor$，在该表达式中必须把尺度差异计入 $E$。恢复 $M$ 要求明确验证该误差小于 $p^{\varepsilon-r}/2$，并且整个同态电路的中间密文都必须有效。

\subsubsection{CKKS 自举}
\label{subsubsec:rns-ckks-bootstrapping}

采用 \autoref{subsubsec:ckks-bootstrapping-summary} 的 CKKS 流程：模数提升、系数到槽位变换、近似模约减，以及槽位到系数变换。稀疏自举还需要 \autoref{subsubsec:ckks-bootstrapping-time-reduction} 中说明的部分求和操作。该基本 CKKS 构造没有单独显式计算秘密密钥内积的同态步骤，也没有普遍适用的末尾重新解释步骤。

若近似 RNS 提升使分量分别增加 $q_0U_A,q_0U_B$，提升后的相位为
\[
\Delta M+E+q_0(K+U_B+U_AS).
\]
逐系数的额外提升界会经过秘密多项式放大。即使 $U_A,U_B$ 的每个系数都属于 $\{-1,0,1\}$，相应相位倍数也未必如此。周期模约减电路的逼近区间必须覆盖全部溢出，误差预算必须包含线性变换、重缩放、密钥切换、多项式逼近，以及保留的输入误差。仅有周期性，并不能使有限的同态多项式成为精确模约减函数。


\subsubsection{BGV 自举}
\label{subsubsec:rns-bgv-bootstrapping}

遵循 \autoref{subsec:bgv-bootstrapping} 的 BGV 构造。首先进行有效的模数切换，使 $\hat q\equiv1\pmod{p^\varepsilon}$，再计算中心化公开系数 $A'=[p^{\varepsilon-r}A]_{\hat q}$、$B'=[p^{\varepsilon-r}B]_{\hat q}$。该特殊模数不必是当前 RNS 基的成员：精确公开重构配合同余保持的模数切换，是一种有效实现途径。

使用明文模数 $p^\varepsilon$、足够大工作密文模数 $Q$ 下的自举密钥，同态计算 $B'+A'S$。RNS 表示的是 $Q$，不要求明文模数等于某个 RNS 因子。随后实现适当的系数与槽位变换，以及采用平衡舍入的数位提取，同时验证 $K'$ 和所有同态中间相位的界。

数位提取期间，每次具有正确性依据的对 $p$ 精确除法，都按 $p^j\rightarrow p^{j-1}$ 改变明文模数；经过 $\varepsilon-r$ 次后，明文模数变为 $p^r$。不能再附加一次乘以 $p^{-(\varepsilon-r)}$ 的操作，否则会对消息重复相除。此外，当 $\hat q$ 不整除 $Q$ 时，基扩展原语 $q\rightarrow qb$ 不能直接描述任意的 $\hat q\rightarrow Q$ 转换。采用这类提升的其他自举构造需要独立的正确性证明，仅有 RNS 表示并不能建立该证明。


\subsubsection{RNS 运算对噪声的影响}
\label{subsubsec:rns-noise}

应跟踪每个原语的实际作用。\textsf{FastBConv} 选择与原值相差源模数有界倍数的提升；\textsf{SmallMont} 限制该倍数，但通常不将其消除。\textsf{ModSwitch\textsubscript{RNS}} 在除法后将基转换偏差变为有界加性误差。\textsf{FastBConvEx} 在范围界成立时是精确的，不需要 \textsf{SmallMont} 来去除误差。\textsf{ModDrop\textsubscript{RNS}} 是精确的剩余类约减，但新相位或尺度条件不成立时，可能使密文失效。

对于分量扰动 $\delta_A,\delta_B$，引起的相位扰动为 $\delta_B+\delta_AS$，其系数界是 $\|\delta_B\|_\infty+\|\delta_A\|_\infty\|S\|_1$。高阶密文还会相应涉及 $S$ 的幂。这些扰动最终被抵消、近似消除，还是保留为解密误差，取决于使用它们的方案和运算，不能用单一的笼统噪声结论涵盖全部 RNS 原语。

\subsubsection{RNS 基本操作的 Python 源代码}
\label{subsubsec:rns-source-code}

我们提供了一个 \href{https://github.com/wang2000star/fhe-textbook-cn/blob/chinese-edition/source%20code/rns.py}{\underline{Python 脚本}}，实现以下 RNS 基本操作示例：\textsf{FastBConv}、
\textsf{ModRaise\textsubscript{RNS}}、
\textsf{ModDrop\textsubscript{RNS}}、
\textsf{ModSwitch\textsubscript{RNS}}、
\textsf{SmallMont} 和
\textsf{FastBConvEx}。





